\chapter{Economic Assessment}\label{sec:EconomicAssessment}

This chapter presents an economic analysis of the resources needed in carrying 
out this Master's thesis, considering the student's working hours, operational 
expenses, and experimental materials. All monetary values 
are expressed in euros.

\section{Labour Cost}

This thesis corresponds to 30 \ac{ECTS} credits. As established in Article 4.5 of 
Spanish Royal Decree~\cite{boe2003rd1125}, one \ac{ECTS} credit 
represents between 25 and 30 hours of student 
workload; the lower bound of 25 hours/\ac{ECTS} is adopted here, giving a total 
estimated workload of $30 \times 25 = 750$ hours. Following the assessment guidelines 
provided for the elaboration of this section, 
the student's time is valued at a nominal rate of $15\,Euro/\text{h}$, even though 
no wage was paid. Table~\ref{tab:LabourCost} distributes the 
750 hours across the main project phases, proportionally to the duration of each 
phase in the project timeline presented in Figure~\ref{fig:GanttChart}.

\begin{table}[H]
    \centering
    \begin{tabular}{lrr}
        \toprule
        \textbf{Phase} & \textbf{Hours} & \textbf{Cost [Euro]} \\
        \midrule
        Literature review \& background research                    & 187.5 & 2{,}812.50 \\
        MATLAB software \ac{EKF} framework adaptation               & 281.3 & 4{,}218.75 \\
        Methodology drafting \& \ac{SOH} algorithm development       &  93.8 & 1{,}406.25 \\
        Experimental testing \& associated results writing           & 125.0 & 1{,}875.00 \\
        Final \ac{EKF} integration, results \& thesis completion     &  62.5 &   937.50 \\
        \midrule
        \textbf{Total} & \textbf{750.0} & \textbf{11{,}250.00} \\
        \bottomrule
    \end{tabular}
    \caption{Labour cost breakdown by project phase, at $15\,Euro/\text{h}$.}
    \label{tab:LabourCost}
\end{table}

\section{Operational Expenses}

The cell was recharged approximately 20 times over the course of the project 
before most test sessions, drawing 
approximately $111.4\,\text{Wh}$ per recharge, as measured directly during the 
charge test (Equation~\ref{eq:EChargeDiscrete}). The power analyser and current 
clamp were estimated to be in active use for approximately 40 hours in total, at 
an assumed power of $20\,\text{W}$.

\begin{table}[H]
    \centering
    \begin{tabular}{lr}
        \toprule
        \textbf{Source} & \textbf{Energy [kWh]} \\
        \midrule
        Cell recharging (20 cycles $\times$ 111.4~Wh) & 2.228 \\
        Laboratory instruments (40~h $\times$ 20~W)   & 0.800 \\
        \midrule
        \textbf{Total} & \textbf{3.028} \\
        \bottomrule
    \end{tabular}
    \caption{Estimated electricity consumption.}
    \label{tab:Electricity}
\end{table}

Using the 2026 annual average regulated electricity price of 
$0.128\,Euro/\text{kWh}$~\cite{PVPC2026}, the total electricity cost was:

\begin{equation}
    3.028\,\text{kWh} \times 0.128\,Euro/\text{kWh} = 0.387\,Euro
\end{equation}

\section{Experimental Materials}

The cell characterisation setup required nine high-power resistors, used as the 
discharge load bank for all discharge tests, at an approximate 
unit cost of $4.93\,Euro$ each, provided by the eRacing team:

\begin{equation}
    9 \times 4.93\,Euro = 44.37\,Euro
\end{equation}

The test setup also included a Texas Instruments microcontroller and six power 
inductors, which were provided by the eRacing team as part of the existing test bench, 
and two Welion SHP270--30 cells, also provided by the team. The individual purchase 
costs of these items were not available at the time of writing and are therefore 
not itemised in this work.

\section{Software Licensing}

MATLAB was used throughout this work for algorithm implementation and 
simulation. Although the licence used was provided by the university, 
its economic value 
is estimated using the listed MathWorks Student annual subscription for 2026 of 
$85\,Euro/\text{year}$~\cite{matlabpricing2026}, 
prorated for approximately $5.3$ months during which the software was actively used for 
this work:

\begin{equation}
    85\,Euro/\text{year} \times \frac{5.3}{12} = 37.54\,Euro
\end{equation}

\section{Total Cost}

Table~\ref{tab:TotalCost} summarises the total estimated cost of this thesis.

\begin{table}[H]
    \centering
    \begin{tabular}{lr}
        \toprule
        \textbf{Item} & \textbf{Cost [Euro]} \\
        \midrule
        Student labour (750~h $\times$ 15~Euro/h)     & 11{,}250.00 \\
        Electricity                                 &      0.39 \\
        Experimental materials (resistors)          &     44.37 \\
        MATLAB licence (prorated)              &     37.54 \\
        \midrule
        \textbf{Total (excl.\ VAT)}                 & \textbf{11{,}332.30} \\
        \bottomrule
    \end{tabular}
    \caption{Total estimated cost of the thesis.}
    \label{tab:TotalCost}
\end{table}

The total cost of the work amounts to $11{,}332.30\,Euro$.

\section{Economic Assessment Discussion}

The labour cost accounts for approximately $98\%$ of the 
estimated budget. The majority of 
the effort was intellectual and computational (literature review, algorithm 
development, and thesis writing) rather than reliant on equipment or 
materials, since the experimental setup used cells, instrumentation, and a test 
bench already available to the eRacing team and the university, at minimal 
marginal cost to this specific project.

% At $15\,€/\text{h}$, the total labour cost of $11{,}250\,€$ for 750 hours is 
% considerably below what an equivalent number of hours would cost at a typical 
% junior engineer's billing rate in Spain (commonly in the range of 
% $25$--$40\,€/\text{h}$), suggesting that this figure represents a conservative 
% lower bound on the real economic value of the work performed. On the other hand, 
% the reuse of existing laboratory infrastructure (power analyser, current clamp, 
% test bench components, and battery cells already owned by the eRacing team) kept 
% the marginal experimental cost extremely low, at under $60\,€$ in total. This 
% highlights one of the practical advantages of the algebraic parameter 
% identification method used in this work: because it requires substantially 
% shorter tests than conventional characterisation approaches, the associated 
% equipment usage time, and therefore its amortised cost, is correspondingly 
% reduced.

% Overall, the total cost obtained is reasonable for a Master's thesis of this scope 
% and is consistent with projects of a similar academic and technical nature, in 
% which most of the value generated stems from the researcher's own time rather 
% than from capital expenditure.